\section{The genuine logarithmic all-degree base-ideal supremum}

Fix the original test configuration, its original invertible total-space
module $M$, its original action linearization, and its fibre-one
identification with the original line $L$. These same original data
are retained in all statements below.

\begin{definition}\label{def:native-tc-original-normalized-log-ideal}
\lean{EnvelopeProperty.NativeTestConfiguration.TestConfiguration.gaussPowerBaseIdealLogNormalized}
\leanok
For the actual base ideal $\mathfrak a_m$ generated by all original
global sections of $M^{\otimes m}$, define at the original Gauss point
\[
 u_m(x)=\frac1m\log\lVert\mathfrak a_m\rVert_{r,x}.
\]
Values lie in the extended reals, and a zero ideal norm gives negative
infinity in every positive degree. The original proper total model
and bounded polynomial Gauss base define the genuine ideal norm.
\uses{def:native-test-configuration-gauss-power-norms}
\end{definition}
\begin{proof}
The original test-configuration projection is proper. Its original
polynomial Gauss valuation is bounded by one for $0<r\leq1$.
Every actual tensor power of $M$ is invertible, so its entire original
global-section base ideal has its genuine ideal norm in $[0,1]$.
Apply the extended logarithm and multiply by the inverse real degree.
\end{proof}

\begin{lemma}\label{lem:native-tc-original-log-root-identity}
\lean{EnvelopeProperty.NativeTestConfiguration.TestConfiguration.gaussPowerBaseIdealLogNormalized_eq_rootLog}
\leanok
The actual normalized logarithm equals the extended logarithm of the
previously constructed real normalized root, embedded in the extended
nonnegative reals. The equality includes zero ideal norms and degree zero.
\uses{def:native-tc-original-normalized-log-ideal,thm:native-test-configuration-actual-gauss-norm-supremum}
\end{lemma}
\begin{proof}
The genuine ideal norm is at most one and hence finite, so embedding
its real value back into the extended nonnegative reals recovers
that same actual norm. Taking a nonnegative real power commutes with
this embedding. The extended logarithm of a real power is that real
exponent times the extended logarithm. The exponent is $1/m$, which
equals the inverse real degree. At degree zero both conventions give
the logarithm of one, namely zero.
\end{proof}

\begin{theorem}\label{thm:native-tc-original-log-limit-and-supremum}
\lean{EnvelopeProperty.NativeTestConfiguration.TestConfiguration.exists_gauss_log_normalized_supremum}
\leanok
If the original fibre-one line is ample, then at every actual original
point and $0<r\leq1$ there is a finite real $a\leq0$ such that
$u_m(x)$ converges to $a$ and
\[
 a=\sup_{m\geq1}u_m(x).
\]
The supremum ranges over every positive original degree, including
degrees with zero base-ideal norm. No convergence or boundedness
certificate is supplied.
\uses{lem:native-tc-original-log-root-identity,thm:native-test-configuration-actual-gauss-norm-supremum}
\end{theorem}
\begin{proof}
The genuine original ampleness theorem supplies a positive real
$b\leq1$ which is the limit and least upper bound of all positive
normalized roots. Put $a=\log b$; positivity makes it finite and
$b\leq1$ gives $a\leq0$. Embedding nonnegative real roots into the
extended nonnegative reals is continuous. The extended logarithm is
a homeomorphism, so the preceding identity transports the genuine
root convergence to the asserted logarithmic convergence. Monotonicity
of the embedding and logarithm transports every positive-degree upper
bound. If $c$ is any upper bound of all positive-degree logarithms,
then all sufficiently late terms are at most $c$. Taking the proven
limit gives $a\leq c$ in the extended reals. This proves the full
least-upper-bound assertion without assuming positivity of early
base-ideal norms.
\end{proof}
